\chapter{Complexity and acceleration of the backward flow}\label{chap:cx}\chaptermark{Complexity and acceleration}

Chapter~\ref{chap:lo} turned the backward Cauchy flow into a positive linear inverse problem and solved it by the EM algorithm of Vardi and Lee. This chapter asks three questions about that computation. What does it cost, and in particular what does log concave sampling in the sense of Applegate, Kannan and Polson say about it? Can the EM iteration be accelerated, by Nesterov momentum or otherwise, and what does the accelerated iteration compute? And does the representation of the Cauchy kernel as a scale mixture of normals, the oldest tool in the Bayesian robustness literature, help?

The chapter begins by stating the deconvolution problem precisely, including the matrix that the computations invert. It goes on to pictures of the Thorin measures at levels $1$, $7/8$, $3/4$ and $1/2$, a theorem on what their positivity means, an integer valued mass deficit that decides positivity up to any finite height, a level mass function whose jumps sit at the real parts of off line zeros, a certification in ball arithmetic of these statements up to height $1419.9$, past the thousandth zero, and a frank account of the limits of the computation. One theorem closes the first part, in three forms: the level one Thorin density has its global minimum at the origin, unconditionally (Proposition~\ref{bp:prop:v1min}); given the quasi Riemann hypothesis, announced and not refereed, the same holds at every level in $[7/8,1]$; and in general, for every $\theta\in(\tfrac12,1]$, the zero free half plane $\Re s>\theta$ is equivalent to the global minimum of $u_{\theta'}$ at the origin for all $\theta'\in[\theta,1]$, with the Riemann hypothesis as the case $\theta\downarrow\tfrac12$ (Theorem~\ref{thm:dbeta:minall}). An a posteriori bound certifies the window masses and centroids of the backward flow from its residual, in ball arithmetic. A smoothed prime sum evaluates $u_1$ at any height to a fixed accuracy in time polynomial in the logarithm of the height, given that the zeros near that height lie on the critical line (verified below $3\times10^{12}$; at every height conditionally on the announced half plane $\Re s>7/8$), and a nonuniform FFT gives a million heights in a second (the author's timing, not rerun); the same explicit formula, read backwards, is the heuristic reason why $\mathfrak U_{1/2}$ (equal to $2\mathfrak U_{1/2}$ of Theorem~\ref{lv:thm:semigroup}, see the notation below) cannot be computed efficiently by deconvolution, and how the conjecture of Hilbert and P\'olya reads in these terms.

Then the answers. The cost splits into three layers. Evaluating the level one measure needs $\zeta$ on the line $\Re s=1$, for which the best algorithms are polynomial in the height $T$ and not in $\log T$; this is the bottleneck and the Thorin formulation does nothing to it. The inversion, given the data, is a convex minimisation, and its Bayesian version is a log concave sampling problem, to which the theorem of Applegate, Kannan and Polson applies with a ratio of radii that is not bounded here; the ill posedness appears as a precision requirement that is linear in the resolution for the linear inverse, hence polynomial in the bits requested. Verification is Turing's method and costs $T^{1+\varepsilon}$. The EM iteration is a majorise minimise scheme in the Kullback Leibler geometry; Nesterov momentum applied in the logarithmic domain keeps positivity, converges an order of magnitude faster at $\theta'=1/2$ in the runs, and in every run reached the same object, the reverse information projection of the data onto the cone of Cauchy blurred positive measures, which is the Thorin measure at level $\theta'$ when it exists (Proposition~\ref{cx:prop:what}). The scale mixture of normals turns the Cauchy blur into a heat flow at a random time with the first passage law of Brownian motion, gives a conditionally Gaussian data augmentation sampler with explicit inverse gamma updates, and identifies the Thorin measure at level $\theta$ with the subordinated heat trace of the zeros; it changes nothing in the complexity classification, and nothing in the arithmetic.

The chapter develops a source chapter on complexity and acceleration of the backward flow by Nick Polson, together with the computations done for this book. Its numbers were recomputed or checked here by scripts written from the mathematics alone (the scripts of the source were not used): \texttt{code/cx\_em.py} for the EM tables of Section~\ref{cx:sec:accel}, \texttt{code/cx\_u\_certify.py} for Theorem~\ref{cx:thm:certified}, \texttt{code/cx\_min.py} for the global minimum certificates and \texttt{code/cx\_slope\_X.py} for Proposition~\ref{cx:prop:slope}, with outputs \texttt{data/cx\_*.json}, listed in \texttt{code/README.md}. Every number printed is one of those recomputations unless it is marked as the author's, and where a recomputation differed from the source the recomputed value is printed. The a posteriori certificate (Theorem~\ref{cx:thm:apostcert}) was computed by the author and was not rerun for this book. Several parts of the source are in the book already and are cited, not restated: the semigroup of Cauchy convolutions and the Thorin measure (Theorem~\ref{lv:thm:semigroup}, Theorem~\ref{lv:thm:thorin}), the existence criterion for the deconvolution (Theorem~\ref{tf:thm:deconv}), the density, the backward multiplier and the EM set up of Chapter~\ref{chap:lo}, the global minimum of the level one density and its equivalent forms (Proposition~\ref{bp:prop:v1min}, Theorem~\ref{thm:dbeta:minall}) and the zero counts (HSW22 \cite{HSW22}, in the form used in Chapter~\ref{chap:basepoint}). What this chapter adds is the theory of EM for the backward flow with its repaired proofs, the acceleration experiments, the analysis of the von Neumann and Ulam estimator, the mass deficit identity, and the certificates of the positivity of $u_{3/4}$ and $u_{7/8}$ on $[0,1420]$.

The notation is the book's, with the following warnings. A level $\theta$ is the real part of a point $s=\theta+iv$, and $u_\theta(v)=\pi^{-1}\Re(\xi'/\xi)(\theta+iv)$ is the density of Chapters~\ref{chap:lv} to~\ref{chap:cd}; in the notation of Chapter~\ref{chap:tf} it is $\rho_{\theta-1/2}(v)$, the ordinate being $v$ here. The Cauchy density is $P_a(x)=(a/\pi)/(a^2+x^2)$. The symmetric Thorin measure $\mathfrak U_\theta$ of equation~\eqref{lo:eq:Utheta} has density $\sum_\rho P_{\theta-\sigma_\rho}(v-\gamma_\rho)$ with no factor $\tfrac12$, so $\mathfrak U_\theta=2\widetilde U_\theta$ with the $\widetilde U_\theta$ of Theorem~\ref{lv:thm:semigroup}. The source's Chapters 2 and 3 are Chapters~\ref{chap:lv} and~\ref{chap:lo} of the book. The word certified is used only for a computation in ball arithmetic whose sentence says what the certificate covers. The announced zero free half plane $\Re s>7/8$ (OpenAI \cite{OAI26a,OAI26b}, with a formalisation in Lean, and Liu \cite{Liu26a}; announced in October 2026 and not yet refereed) is never used as a theorem: a result that uses it is stated as conditional on it.

Status words follow the rest of the book: proved, computed, conditional, heuristic or open. The results are these.
\begin{enumerate}
\item The Thorin measure at level $\theta'$ is the solution of a positive deconvolution; the matrix $A$ of its discretisation, with cell densities as unknowns, has the properties (a) to (e) of Section~\ref{cx:sec:problem} (proved), including the kernel of dimension exactly $2(N-M)$ when the unknown window is larger. The level one data must be supplied to absolute precision $e^{-(1-\theta')\Omega}$ for the linear inverse (Proposition~\ref{cx:prop:precision}, proved). Log concavity of the Poisson posterior is proved (Theorem~\ref{cx:thm:logconcave}); the sampling clause is a cited theorem and is conditional on a bound for the ratio of radii that is not proved here.
\item EM is a majorise minimise scheme with the rate $D(x^*\|x^{(0)})/k$ in the consistent case (Theorem~\ref{cx:thm:emrate}, proved, with the repaired step) and local linear convergence in the interior (Theorem~\ref{cx:thm:linear}, proved; its injectivity hypothesis fails in the runs). What EM computes: under three growth conditions its limits converge vaguely to $\mathfrak U_{\theta'}$ when $\zeta$ has no zeros in $\Re s>\theta'$, and the residual stays bounded away from zero when it has one (Proposition~\ref{cx:prop:what}, proved).
\item Log domain Nesterov momentum with restart and SQUAREM are an order of magnitude faster than EM at $\theta'=1/2$ on a window; on the full range ($202$ zeros below height $400$) and on the first thousand zeros ($994$ of $995$ zeros below $1415$ found as peaks) they locate the zeros to about $10^{-2}$ (Section~\ref{cx:sec:accel}, computed; the rates $O(1/k)$ and $O(1/k^2)$ and the identification of the limit are observed, not proved).
\item The von Neumann and Ulam estimator for the backward flow has infinite variance whenever $h<0.5221\,\pi(1-\theta')$ (Proposition~\ref{cx:prop:vnvar}, proved); the positive reverse walk is the EM step in expectation. The Cauchy kernel is a first passage mixture of Gaussians, giving the subordinated heat trace of the zeros and an inverse gamma data augmentation sampler (Proposition~\ref{cx:prop:subord}, proved; the sampler comparison is computed).
\item The level $7/8$ Thorin measure is computed by three routes (computed); its identification as a Thorin measure is conditional on the announced half plane. The positivity theorem (Theorem~\ref{cx:thm:positivity}), the mass deficit identity (Proposition~\ref{cx:prop:deficit}) and the level mass function (Corollary~\ref{cx:cor:levelmass}) are proved; the deficit tables, the $U_{3/4}$ and $U_{1/2}$ tables and Turing's check in Thorin bookkeeping are computed. Decidability by level is proved for finite heights (Section~\ref{cx:sec:limits}).
\item Certified positivity (Theorem~\ref{cx:thm:certified}, computed in ball arithmetic at two precisions, $96$ and $160$ bits): the densities $u_{3/4}$ and $u_{7/8}$ are positive on $[0,1420]$, every certified lower bound exceeding $4\times10^{-5}$ in the recomputation; the argument change of $\xi$ gives exactly $1000$ zeros below height $1419.9$, all on the critical line and simple, which is known (Platt and Trudgian \cite{PT21}); the positivity of $u_{3/4}$ and $u_{7/8}$ on this range is what the certificate adds.
\item The global minimum of the Thorin density is at the origin: at level $1$ this is Proposition~\ref{bp:prop:v1min} (the proof of Theorem~\ref{cx:thm:globalmin}(i) is an independent computed one); at level $7/8$ it holds conditionally on the announced half plane and is computed, with least certified bound on $[13.5,1420]$ equal to $0.0360$ at $\theta=1$ and $0.0272$ at $\theta=7/8$, both at $v=17.49$ (Theorem~\ref{cx:thm:globalmin}); the Riemann hypothesis form is Theorem~\ref{thm:dbeta:minall}(vi) (Theorem~\ref{cx:thm:rhmin}). The uniform law and the lower bound for the excess, $W\ge w^*$ with its brackets (Theorems~\ref{cx:thm:uniform} and~\ref{cx:thm:lbl}), are computed, and conditional on the announced half plane at level $7/8$.
\item The slope of the level one density has certified infima near $0.1758$ to $0.1793$ (Proposition~\ref{cx:prop:slope}, computed), and under the Riemann hypothesis the limiting bracket $X$ has its minimum in the first gap at $v^*=17.4979\pm0.0001$, with $0.1838985\le X(v^*)\le0.1838988$ (computed, conditional on RH); that this is its infimum over $13.5\le|v|\le1418.9$ is the author's computation, not rerun.
\item The a posteriori bound for the window masses and centroids of the backward flow, from its residual, is proved (Proposition~\ref{cx:prop:apost}); the certificate of Theorem~\ref{cx:thm:apostcert} is computed by the author and not rerun here, and is conditional on the absence of zeros to the right of $\theta'$ (the Riemann hypothesis at $\theta'=\tfrac12$); it certifies Gaussian weighted masses and centroids, and single zeros only through Corollary~\ref{cx:cor:zeroloc} (proved).
\item The smoothed prime sum for $u_1$ has an exact identity (proved) and an error bound proved under the local hypothesis on the zeros, with the pole term that the source had omitted (Theorem~\ref{cx:thm:fastu1}); conditionally on the announced half plane it holds at every height with $X$ of order $(\log v/\varepsilon)^8$; the nonuniform FFT experiment is computed (by the author, not rerun).
\item The explicit formula read backwards, and the efficiency of the backward flow (Section~\ref{cx:sec:halfeff}), and the reading of the Hilbert and P\'olya conjecture as a resolvent trace (Section~\ref{cx:sec:vonneumann}), are heuristic. The open problems close the chapter.
\end{enumerate}
